\subsection*{Question 2.7 - Transfer functions}

We use the linearized system of Question 2.6, written element by element, with $\tilde x = (\tilde x, \tilde y, \tilde z, \tilde\phi, \tilde\theta, \tilde\psi, \tilde u, \tilde v, \tilde w, \tilde p, \tilde q, \tilde r)^T$ and $\tilde u = (\tilde T, \tilde n_x, \tilde n_y, \tilde n_z)^T$.
(Here $\tilde u$ and $\tilde v$ in the state denote the first two components of the velocity error, not the input vector or the whole velocity.)
Since the system is linear and time invariant, each transfer function is obtained by taking the Laplace transform of the relevant equations with zero initial conditions.
We assume $J = \mathrm{diag}(J_{xx}, J_{yy}, J_{zz})$, so $J^{-1}$ is diagonal and the dynamics of the three rotation axes are decoupled.

\paragraph{Step 1 - Attitude channels.}
The last two block rows of the system are $\dot{\tilde\lambda} = \tilde\omega$ and $\dot{\tilde\omega} = J^{-1}\tilde n$.
Componentwise they read
\begin{equation*}
\ddot{\tilde\phi} = \frac{\tilde n_x}{J_{xx}}, \qquad
\ddot{\tilde\theta} = \frac{\tilde n_y}{J_{yy}}, \qquad
\ddot{\tilde\psi} = \frac{\tilde n_z}{J_{zz}} .
\end{equation*}
In the Laplace domain every double derivative becomes a multiplication by $s^2$, so
\begin{equation*}
G_\phi(s) = \frac{\Phi(s)}{N_x(s)} = \frac{1}{J_{xx}\, s^2}, \qquad
G_\theta(s) = \frac{\Theta(s)}{N_y(s)} = \frac{1}{J_{yy}\, s^2}, \qquad
G_\psi(s) = \frac{\Psi(s)}{N_z(s)} = \frac{1}{J_{zz}\, s^2}.
\end{equation*}
Each attitude angle behaves as a double integrator driven by its own torque.

\paragraph{Step 2 - Vertical channel.}
The third components of $\dot{\tilde p} = \tilde v$ and of $\dot{\tilde v} = A_{v\lambda}\tilde\lambda + b_T\tilde T$ are $\dot{\tilde z} = \tilde w$ and $\dot{\tilde w} = -\tilde T/m$.
(The third row of $A_{v\lambda}$ is zero, so the vertical channel does not depend on the attitude.)
Combining them, $\ddot{\tilde z} = -\tilde T / m$, and
\begin{equation*}
G_z(s) = \frac{Z(s)}{T(s)} = -\frac{1}{m\, s^2}.
\end{equation*}
The minus sign appears because $z$ points down: more thrust makes the vehicle go up, that is, towards negative $z$.

\paragraph{Step 3 - Horizontal channels.}
The first two components of the velocity equations are $\dot{\tilde u} = -g\,\tilde\theta$ and $\dot{\tilde v} = g\,\tilde\phi$.
Together with $\dot{\tilde x} = \tilde u$ and $\dot{\tilde y} = \tilde v$ they give
\begin{equation*}
\ddot{\tilde x} = -g\,\tilde\theta, \qquad \ddot{\tilde y} = g\,\tilde\phi .
\end{equation*}
In the Laplace domain, $X(s) = -\frac{g}{s^2}\,\Theta(s)$ and $Y(s) = \frac{g}{s^2}\,\Phi(s)$.
Substituting $\Theta(s) = G_\theta(s) N_y(s)$ and $\Phi(s) = G_\phi(s) N_x(s)$ from Step 1,
\begin{equation*}
G_x(s) = \frac{X(s)}{N_y(s)} = -\frac{g}{J_{yy}\, s^4}, \qquad
G_y(s) = \frac{Y(s)}{N_x(s)} = \frac{g}{J_{xx}\, s^4}.
\end{equation*}
Note that $\tilde x$, $\tilde y$ and $\tilde z$ are the components of the position error $\tilde p$ defined in Question 2.6, that is, $R^T(p - p_0)$, and not of the inertial position error.
The horizontal position responds to the torque with four integrators: two from the torque to the angle, and two from the angle (through the tilted thrust) to the position.
